\section*{Chapter \ref{ch:g12:graph-algorithms} --- Graph Traversals and Shortest Paths}

\begin{solution}{exo:g12:graph-algorithms:1}
D, B, A, C, E, F. From D, the first neighbour is B; from B, A; from A, C (B is visited); from C, E (A and D are
visited); from E, F.
\end{solution}

\begin{solution}{exo:g12:graph-algorithms:2}
Layer $0$: D. Layer $1$: B, C, F. Layer $2$: A (found from B), E (found from C). The order is D, B, C, F, A, E.
\end{solution}

\begin{solution}{exo:g12:graph-algorithms:3}
Rows and columns in the order $1$ to $4$:
\[
\begin{pmatrix} 0 & 1 & 1 & 1 \\ 1 & 0 & 1 & 0 \\ 1 & 1 & 0 & 1 \\ 1 & 0 & 1 & 0 \end{pmatrix}
\]
and the lists \code{{1: [2, 3, 4], 2: [1, 3], 3: [1, 2, 4], 4: [1, 3]}}.
\end{solution}

\begin{solution}{exo:g12:graph-algorithms:4}
E $0$; C and F $1$; A and D $2$; B $3$.
\end{solution}

\begin{solution}{exo:g12:graph-algorithms:5}
\smallskip\noindent\hspace*{2em}\code{def count_edges(adj):}

\noindent\hspace*{4em}\code{return sum(len(ns) for ns in adj.values()) // 2}

\smallskip\noindent Each edge appears in the lists of both its ends: $14 / 2 = 7$ for \code{G}.
\end{solution}

\begin{solution}{exo:g12:graph-algorithms:6}
$16$: one call per square, since a square once marked is never entered again. The student's version unmarks a square
when it returns, so the square can be entered again along another path: it walks every path from the corner that
never crosses itself. Their number grows very fast with the size of the grid: $7$, $79$ and $2111$ calls for the
grids of $2 \times 2$, $3 \times 3$ and $4 \times 4$ squares.
\end{solution}

\begin{solution}{exo:g12:graph-algorithms:7}
\smallskip\noindent\hspace*{2em}\code{def has_path(g, a, b):}

\noindent\hspace*{4em}\code{seen, stack = {a}, [a]}

\noindent\hspace*{4em}\code{while stack:}

\noindent\hspace*{6em}\code{v = stack.pop()}

\noindent\hspace*{6em}\code{if v == b:}

\noindent\hspace*{8em}\code{return True}

\noindent\hspace*{6em}\code{for w in g[v]:}

\noindent\hspace*{8em}\code{if w not in seen:}

\noindent\hspace*{10em}\code{seen.add(w)}

\noindent\hspace*{10em}\code{stack.append(w)}

\noindent\hspace*{4em}\code{return False}
\end{solution}

\begin{solution}{exo:g12:graph-algorithms:8}
Finish A ($0$): S $4$, B $2$, C $5$. Finish B ($2$): S becomes $3$, D $10$. Finish S ($3$): nothing changes. Finish C
($5$): D becomes $6$, T $8$. Finish D ($6$), then T ($8$). Distances from A: S $3$, B $2$, C $5$, D $6$, T $8$.
\end{solution}

\begin{solution}{exo:g12:graph-algorithms:9}
After S, it finishes B, the last found, at $1$: correct by luck. A drops to $3$ and D gets $9$; it finishes D, the
last found, although A, at $3$, is closer. From D, C gets $10$ and T $15$. It then finishes T at $15$, while the route
S, B, A, C, T of length $11$ is not yet known, and C at $10$ instead of $8$. Finishing a vertex declares its distance
final, which is only safe for the closest unfinished vertex.
\end{solution}

\begin{solution}{exo:g12:graph-algorithms:10}
The lists hold one entry per direction of each road: $52 \times 2 + 1 = 105$ entries, plus the $30$ places. The matrix
has $30 \times 30 = 900$ cells, of which $795$ are $0$.
\end{solution}

\begin{solution}{exo:g12:graph-algorithms:11}
\smallskip\noindent\hspace*{2em}\code{def bfs_layers(g, s):}

\noindent\hspace*{4em}\code{layers, seen = [[s]], {s}}

\noindent\hspace*{4em}\code{while True:}

\noindent\hspace*{6em}\code{nxt = []}

\noindent\hspace*{6em}\code{for v in layers[-1]:}

\noindent\hspace*{8em}\code{for w in g[v]:}

\noindent\hspace*{10em}\code{if w not in seen:}

\noindent\hspace*{12em}\code{seen.add(w)}

\noindent\hspace*{12em}\code{nxt.append(w)}

\noindent\hspace*{6em}\code{if not nxt:}

\noindent\hspace*{8em}\code{return layers}

\noindent\hspace*{6em}\code{layers.append(nxt)}

\smallskip\noindent No queue is needed: each layer is built from the one before. For \code{G} from D:
\texttt{[[\textquotesingle{}D\textquotesingle{}], [\textquotesingle{}B\textquotesingle{}, \textquotesingle{}C\textquotesingle{}, \textquotesingle{}F\textquotesingle{}], [\textquotesingle{}A\textquotesingle{}, \textquotesingle{}E\textquotesingle{}]]}.
\end{solution}

\begin{solution}{exo:g12:graph-algorithms:12}
Write $\delta(v)$ for the true distance and $d(v)$ for the number the search records.

First, $d(v) \ge \delta(v)$: following the parents from $v$ gives a real path of $d(v)$ edges.

Next, the queue: we show by induction on the steps that its numbers, from front to back, never decrease and differ by
at most $1$. A step removes the front vertex $v$, of number $d$, and appends vertices numbered $d + 1$; since every
number left was between $d$ and $d + 1$, the property is kept. So the vertices leave the queue in increasing order
of their numbers.

Finally, $d(v) \le \delta(v)$, by induction on $k = \delta(v)$. For $k = 0$, $v = s$. If $\delta(v) = k \ge 1$, let
$u$ be the vertex before $v$ on a shortest path: $\delta(u) = k - 1$, so $d(u) = k - 1$ by induction, and $u$ goes
through the queue. When $u$ is removed, $v$ is numbered, either now as $d(u) + 1 = k$, or earlier by a vertex $x$
removed before $u$, with $d(x) \le d(u)$, so $d(v) = d(x) + 1 \le k$. Hence $d(v) = \delta(v)$.
\end{solution}

\begin{solution}{exo:g12:graph-algorithms:13}
By induction on the steps: assume every vertex finished so far has its true distance, and let $u$ be the next one,
with tentative distance $d(u)$. Tentative distances are lengths of real paths, so $d(u) \ge \delta(u)$. Take any path
$P$ from $s$ to $u$. It starts at a finished vertex and ends at an unfinished one: let $y$ be its first unfinished
vertex, and $x$ the vertex just before it, finished. When $x$ was finished, its edge to $y$ was examined, so
$d(y) \le \delta(x) + w(x, y)$, which is the length of $P$ up to $y$. Then $d(u) \le d(y)$, since $u$ has the smallest
tentative distance; and the rest of $P$, from $y$ to $u$, has weights $\ge 0$. So the length of $P$ is at least
$d(y) \ge d(u)$. This holds for every path: $\delta(u) \ge d(u)$, and $d(u) = \delta(u)$. The step ``rest of $P$ is
$\ge 0$'' is exactly where negative weights break the proof.
\end{solution}

\begin{solution}{exo:g12:graph-algorithms:14}
\smallskip\noindent\hspace*{2em}\code{def topological_order(g):}

\noindent\hspace*{4em}\code{state, order = {v: 0 for v in g}, []}

\noindent\hspace*{4em}\code{def explore(v):}

\noindent\hspace*{6em}\code{state[v] = 1}

\noindent\hspace*{6em}\code{for w in g[v]:}

\noindent\hspace*{8em}\code{if state[w] == 1:}

\noindent\hspace*{10em}\code{raise ValueError("the graph has a cycle")}

\noindent\hspace*{8em}\code{if state[w] == 0:}

\noindent\hspace*{10em}\code{explore(w)}

\noindent\hspace*{6em}\code{state[v] = 2}

\noindent\hspace*{6em}\code{order.append(v)}

\noindent\hspace*{4em}\code{for v in g:}

\noindent\hspace*{6em}\code{if state[v] == 0:}

\noindent\hspace*{8em}\code{explore(v)}

\noindent\hspace*{4em}\code{return order[::-1]}

\smallskip\noindent The three states are the three colours. A vertex is appended once all its \omterm{def:g12:graph-algorithms:adjacency-list}{successors} are, so in
the reversed list it comes before them. The tests compare it with every order of small random \omterm{def:g10:graphs-first-look:graph}{graphs}.
\end{solution}

\begin{solution}{exo:g12:graph-algorithms:15}
Add $3$: S--A weighs $4$, S--B $6$, B--A $0$. \omterm{met:g12:graph-algorithms:dijkstra}{Dijkstra's algorithm} finds S to A $= 4$ by the direct edge, against $6$
through B; in the original \omterm{def:g10:graphs-first-look:graph}{graph} that direct path costs $1$, but the path through B costs $0$. Adding $3$ to each
edge adds $3$ times the number of edges to each path: paths with more edges are penalised more, and the order of the
paths changes.
\end{solution}

\begin{solution}{pb:g12:graph-algorithms:1}
\textbf{1.} The vertices are the $8 \times 14 = 112$ squares; an edge joins two squares side by side with no wall
between them.

\textbf{2.} A tree on $112$ vertices has $111$ edges: $111$ passages. The program counts $119$: of the twelve walls
picked afterwards, eight were real walls, and four were already open passages, which the test \code{if b not in g[a]}
skips.

\textbf{3.} A stack: the carving is a depth-first search, which at each step goes to a random unvisited neighbour, and
backs up from dead ends.

\textbf{4.} The eight extra passages each close a loop in the tree, giving a second route between the squares they
join.

\textbf{5.} $92$ squares out of $112$. Depth-first search follows its neighbours in a fixed order, deep into whatever
branch comes first, with no idea where the exit is: it reaches the exit only when that branch's turn comes.

\textbf{6.} Right first: as long as the wall on the right has an opening, the walker turns into it, so its right hand
never leaves the wall. Ahead if the right is closed, then left, and back only in a dead end.

\textbf{7.} $58$ steps, against $32$: it walks into every dead end that opens on its right side, and back out.

\textbf{8.} The entrance is a corner square, so the walls around it belong to the outer wall and the walls joined to
it, one connected group of walls; the walker keeps its hand on that group, which runs all the way round the maze
and past the exit square, also on the outer wall. The loops create walls standing alone inside the maze, and only a
walker started beside one of them could circle for ever.

\textbf{9.} The wall follower keeps only its square and the direction it faces. A search keeps the set of visited
squares and its stack or queue, a memory that grows with the maze.

\textbf{10.} $32$ steps: $33$ stars, the entrance and the exit included.

\textbf{11.} Each parent is one step closer to the entrance (exercise \ref{exo:g12:graph-algorithms:12}): following
them back from the exit takes $32$ steps, the distance.

\textbf{12.} Not necessarily. Counting them layer by layer, each square has as many shortest routes as the sum of
those of its neighbours one layer closer; the tests run this count on the maze and find $5$ shortest exits.

\textbf{13.} As weights: \code{g[a][b]} is $6$ when the passage between $a$ and $b$ is a door, $1$ otherwise, in both
directions.

\textbf{14.} It counts edges as if they all cost the same; here a door costs six times an open passage.

\textbf{15.} $32$ steps, $57$ seconds: $57 = 32 + 5 \times 5$, so $5$ doors.

\textbf{16.} $34$ steps, $54$ seconds: $54 = 34 + 4 \times 5$, so $4$ doors. Two more steps avoid one door.

\textbf{17.} All weights, $1$ and $6$, are positive: the theorem applies.

\textbf{18.} $32$ seconds: every step costs $1$, and the fastest exit is the one with fewest steps.

\textbf{19.} $n = 112$: at most $112^2 = 12\,544$ comparisons, nothing for a computer.

\textbf{20.} \textbf{The shortest exit of the maze takes $32$ steps, where the hand on the wall takes $58$ (and the
depth-first search visits $92$ squares); with doors, the fastest exit takes $54$ seconds in $34$ steps, against $57$
seconds for the exit with fewest steps.}
\end{solution}
